


\section{Related Work}
\label{sec.related_work}


\begin{table}[t]
\centering
\addtolength{\tabcolsep}{-0.3em}
\resizebox{\columnwidth}{!}{
\begin{tabular}{l c c c c c}
\toprule
Methods & Training & \makecell{Goal-oriented\\Sub-agent} & \makecell{Holistic\\Orchestration} & \makecell{Deg. of \\MAS} & \makecell{Analysis on \\ MAS vs. SAS} \\
\hline
MAS-Zero \cite{Ke2025MASZero} & \xmark & \xmark & \cmark & \xmark & \xmark \\
AFlow \cite{zhang2024aflowautomatingagenticworkflow} & \xmark & \xmark & \xmark & \xmark & \xmark \\
DyLAN \cite{liu2023dynamic} & \xmark & \xmark & \xmark & \xmark & \xmark \\
 \midrule
MAS-GPT \cite{ye2025masgpttrainingllmsbuild} & SFT & \xmark & \cmark & \xmark & \xmark \\
CoRL \cite{jin2025controllingperformancebudgetcentralized} & RL & \xmark & \xmark & \xmark & \xmark \\
Puppet \cite{dang2025multiagentcollaboration} & RL & \xmark & \xmark & \xmark & \xmark \\
xRouter \cite{qian2025xrouter} & RL & \xmark & \xmark & \xmark & \xmark \\
W4S \cite{nie2025weakforstrong} & RL & \xmark & \xmark & \xmark & \xmark \\
ToolOrchestra \cite{su2025toolorchestraelevatingintelligenceefficient} & RL & \cmark & \xmark & \xmark & \xmark \\
\textbf{\ourframework} & RL & \cmark & \cmark & \cmark & \cmark \\
\bottomrule
\end{tabular}
}
\caption{Comparison of popular automatic MAS approaches. 
% \zixuan{or we can have a parato front fig}
}
\label{tab:method_comparison}
\vspace{-2\baselineskip}
\end{table}

\subsection{From Single to Multi-Agent Systems}
While \textit{agenticness} is generally viewed as a spectrum \cite{kapoor2024ai}, it remains foundational to clearly distinguish between Single-Agent System (SAS) and Multi-Agent System (MAS). We define SAS as one that features a single reasoning locus {with one goal}: all perceptions, plans, and actions are carried out within a single sequential control loop {to achieve the goal} governed by one LLM instance. This definition includes systems that employ tool use \citep{yao2023react}, self-reflection \citep{madaan2024self}, or CoT reasoning. In contrast, MAS comprise multiple LLM-backed agents {with their own goals} that interact through structured communication mechanisms, such as message passing, shared memory, or explicitly orchestrated protocols \citep{xi2023risepotentiallargelanguage}. Each agent may maintain its own context, objectives, and tools, and system-level behavior emerges from their coordination {in the form of \emph{collective reasoning}} rather than from a single monolithic reasoning process.

% While MAS is related to other concepts such as modular reasoning or tool composition, its defining characteristic lies in the presence of multiple interacting reasoning entities. \zixuan{there are also some related concepts}

% Early MAS designs were largely manually constructed, with fixed interaction patterns like debate \citep{du2023improvingfactualityreasoning}. More recently, there has been growing interest in \textbf{automatic MAS design}, where the system structure, agent roles, and interactions are generated dynamically. 
Depending on how a sub-agent is defined in different contexts, MAS design also relates to LLM router {\citep{su2025toolorchestraelevatingintelligenceefficient,zhang2025routerr1teachingllmsmultiround}}  (LLM as sub-agent) and Skills \citep{anthropic_agent_skills_2024} (skill as sub-agent). While there is work on manual MAS designs and automataic inference-time orchestration (see \cref{ap:related}), we focus our discussion on \textbf{training-time orchestration}, given its growing popularity and relevance to \ourframework. 
% \ourframework falls into this latter category, focusing on automatic MAS design that enables scalable and principled orchestration of complex sub-agents.

\subsection{Automatic Training-time Orchestration}
% \textbf{Inference-time orchestration.} This line of work typically relies on external signals from validation sets or LLM-based judges to guide adaptation \citep{zhou2025multiagentdesignoptimizingagents,zhang2025multiagentarchitecturesearchagentic,hu2025automated,zhang2024aflowautomatingagenticworkflow,Ke2025MASZero}. However, these methods often exhibit limited adaptability due to unreliable adaptation signals.


% these inference-time methods exhibit limited adaptability, as unreliable validation signals or incorrect guidance from LLM-based judges can misdirect the adaptation process.

% MAS-Zero~\citep{Ke2025MASZero} is an exception, performing inference-time MAS design without validation feedback. 

% Recent work on automatic MAS design typically operates at inference time and relies on a validation set to guide adaptation \citep{zhou2025multiagentdesignoptimizingagents,zhang2025multiagentarchitecturesearchagentic,hu2025automated,zhang2024aflowautomatingagenticworkflow}. MAS-Zero~\citep{Ke2025MASZero} differs by performing inference-time MAS design without validation feedback.
% MASS~\citep{zhou2025multiagentdesignoptimizingagents} uses rejection sampling to select effective agent configurations, and MaAS~\citep{zhang2025multiagentarchitecturesearchagentic} extends MASS with a question-wise masking mechanism to adapt subnetworks. ADAS~\citep{hu2025automated} and AFlow~\citep{zhang2024aflowautomatingagenticworkflow} frame MAS design as a code generation task. ADAS stores and searches historical designs based on validation performance, while AFlow enhances this with Monte Carlo Tree Search. 
% One exception is MAS-Zero \citep{Ke2025MASZero}, which performs fully inference-time automatic MAS design without relying on a validation set. 

% \textbf{Training-time orchestration.} 
This line of work seeks to avoid expensive inference-time adaptation by directly training an orchestrator. Most existing approaches in this category adopt sequential orchestration and optimize it using multi-step RL \citep{su2025toolorchestraelevatingintelligenceefficient,nie2025weakforstrong}.
An exception is MAS-GPT \citep{ye2025masgpttrainingllmsbuild}, which trains an orchestrator without RL; however, it reports lower performance even compared to inference-time adaptation methods \citep{Ke2025MASZero}, suggesting limitations in capturing effective system-level coordination.
To our knowledge, \ourframework is the first approach to perform training-time holistic orchestration, in which the meta-agent is trained to generate a complete multi-agent system in a single decision step. This formulation enables the orchestrator in \ourframework to reason at the {\emph{plan level rather than over execution trajectories}}, allowing for global coordination across sub-agents, avoiding error accumulation from intermediate states, and directly aligning training objectives with end-task performance.
% ... + adavtange \zixuan{avoid duplicaate}
Table~\ref{tab:method_comparison} provides a systematic comparison between \ourframework and representative inference- and training-time MAS design approaches. The comparison shows that \ourframework is the only method that simultaneously supports goal-oriented sub-agents, holistic orchestration, configurable \masness, and controlled experiments aimed at understanding when and how MAS can outperform SAS.
% \zixuan{need to mention}

% Infernce-time, training time...



